\documentclass[10pt]{article}
\usepackage[T1]{fontenc}
\usepackage[utf8]{inputenc}
\usepackage[margin=1in]{geometry}
\usepackage{microtype}
\usepackage{amsmath,amssymb,amsthm}
\usepackage{booktabs,array,longtable}
\usepackage{xurl}
\usepackage[hidelinks]{hyperref}
\usepackage{graphicx}
\setlength{\emergencystretch}{2em}
\newtheorem{proposition}{Proposition}
\newtheorem{definition}{Definition}
\newcommand{\method}{Prometheus}
\newcommand{\R}{\mathbb{R}}
\newcommand{\E}{\mathbb{E}}
\title{Prometheus: Learning and Preserving\\Executable Operations through Monodromy}
\author{Prometheus Research Project}
\date{Working manuscript: 1 November 2026}

% Canonical method: water/method.md. Direction: water/now.md.
% Protocol authority: ALL.md, TODO 5. Source and evidence remain versioned.
% Main text follows the human-specified eight-section paper structure.
% All detailed development records and unvalidated extensions are in the appendix.
\begin{document}
\maketitle
\begin{abstract}
Learning systems can retain individual operations yet lose the ability to call
and compose them when shared representations change. We investigate Prometheus,
a candidate computational architecture that represents executable operations by
monodromy and learns their neural bindings from terminal task feedback. A Block
combines a learned state encoder, shared loop controller, operation calls, and
readout. Operation formation changes a loop's executable class; subsequent
within-class deformation preserves its action through a nonsingular homotopy.
A coupled gradient combines ordinary credit through the neural computation with
loss jumps at class boundaries. Protected update transactions additionally check
old task bindings and soft/native execution consistency, since preserving an
operation alone does not preserve its use. We derive sufficient retention
conditions and quantify workspace requirements for irreversible logical
interfaces. In a bounded integrated diagnostic, the coupled learner acquires a
new operation and calling behavior while retaining its supplied old capability;
a targeted block-coordinate repair resolves a measured feasibility bottleneck
without adding parameters. Historical continual-learning prototypes and matched
ordinary constraints are documented separately. These development results do
not establish topology-specific competitive advantage. We specify a frozen,
five-layer evaluation against Transformers and continual-learning baselines,
including acquisition, forgetting, composition, parameter growth, and complete
resource accounting. The complete general Block and its primary benchmark
remain under development. The central hypothesis is that topology can enlarge
useful shared adaptation while reducing the cost of preserving executable
computational capabilities.
\end{abstract}
\input{sections/01-introduction.tex}
\input{sections/02-related-work.tex}
\input{sections/03-architecture.tex}
\input{sections/04-monodromy-learning.tex}
\input{sections/05-theory.tex}
\input{sections/06-experiments.tex}
\input{sections/07-discussion.tex}
\input{sections/08-conclusion.tex}

\clearpage
\bibliographystyle{plain}
\bibliography{references}
\clearpage
\appendix
\input{appendix/materials.tex}
\end{document}
